\chapter{Nets and compact upper-envelope calculus}
\index{net}\index{Moore--Smith convergence}\index{upper envelope!net calculus}
\label{app:nets}

This appendix gives the point-set and extended-real facts used by the
compactness and envelope arguments.
It is deliberately formulated for arbitrary directed sets.
A sequence may replace a net only after a metrizability, first-countability,
or Choquet-selection argument has been supplied.

The first use is \Cref{chap:topological-tools}.
The distinction between increasing nets and sequences is essential in
\Cref{chap:bj-homogeneity-bdivisors} and in the envelope proof.
Our subnet convention is the Willard tail convention
\cite[Chapter~4]{Wil70}.

\section{Directed sets, tails, and convergence}
\label{sec:appE-directed}

\begin{definition}[Directed preorder]
\label{def:appE-directed}
A directed set is a nonempty set \(A\) with a reflexive transitive relation
\(\leq\) such that every pair has a common upper bound.
Antisymmetry is not required.
A net in a set \(X\) is a map \(a\mapsto x_a\) from \(A\) to \(X\).
The tail beginning at \(a_0\) is
\[
                         \{x_a:a\geq a_0\}.
\tag{E.1}\label{eq:appE-tail}
\]
\end{definition}

A property holds eventually if it holds on one tail, and frequently if it
holds at some index beyond every prescribed index.
Thus \(x_a\) is frequently in \(S\subset X\) if
\[
             \forall a_0\in A\ \exists a\geq a_0
             \quad x_a\in S.
\tag{E.2}\label{eq:appE-frequently}
\]

\begin{definition}[Convergence and cluster point]
\label{def:appE-convergence}
In a topological space \(X\), the net \(x_a\) converges to \(x\) if it is
eventually in every neighborhood of \(x\).
The point \(x\) is a cluster point of the net if the net is frequently in
every neighborhood of \(x\).
\end{definition}

\begin{example}[Neighborhood net]
\label{ex:appE-neighborhood-net}
Direct the neighborhoods of \(x\) by reverse inclusion:
\[
 U\preceq V\quad\Longleftrightarrow\quad V\subseteq U.
\]
Choosing \(x_U\in U\) gives a net converging to \(x\).
This construction, unavailable to sequences without first countability,
is the basic closure test.
\end{example}

\section{Willard and Moore--Smith subnets}
\label{sec:appE-subnets}

\begin{definition}[Willard subnet]
\label{def:appE-willard-subnet}
Let \(x:A\to X\) be a net.
A Willard subnet is \(y_b=x_{\phi(b)}\), indexed by a directed set \(B\),
where \(\phi:B\to A\) is order preserving and eventually in every tail:
\[
 \forall a_0\in A\ \exists b_0\in B\ \forall b\geq b_0,
 \qquad \phi(b)\geq a_0.
\tag{E.3}\label{eq:appE-willard}
\]
The map \(\phi\) need not be injective.
\end{definition}

\begin{definition}[Order-preserving Moore--Smith subnet]
\label{def:appE-ms-subnet}
In this book an order-preserving Moore--Smith subnet means exactly a
Willard subnet: \(\phi:B\to A\) is order preserving and satisfies
\eqref{eq:appE-willard}.
\end{definition}

Terminology varies between topology texts, so the defining formulas, not the
name, control every use in this book.

\begin{lemma}\label{lem:appE-tail-subnet}
If \(x_a\) is eventually in \(S\), every Willard subnet is eventually in
\(S\).
Consequently, if \(x_a\to x\), every subnet still converges to \(x\).
This does not assert uniqueness unless the ambient space is Hausdorff.
\end{lemma}

\begin{proof}
Choose \(a_0\) such that \(x_a\in S\) for \(a\geq a_0\).
Condition \eqref{eq:appE-willard} makes
\(\phi(b)\geq a_0\) eventually.
\end{proof}

\begin{theorem}\label{thm:appE-cluster-subnet}
A point \(x\) is a cluster point of \(x_a\) if and only if some Willard
subnet converges to \(x\).
In fact, one may choose an order-preserving Moore--Smith subnet.
\end{theorem}

\begin{proof}
The reverse implication follows from
\Cref{lem:appE-tail-subnet}: if the original net were eventually
outside a neighborhood \(U\), so would every subnet.

Conversely, let \(B\) consist of triples
\[
 (a,U,c),\qquad a\in A,\quad U\ni x,\quad
 c\geq a,\quad x_c\in U.
\]
Clusterhood makes this set nonempty over every pair \((a,U)\).
Order triples by
\[
 (a,U,c)\preceq(a',U',c')
 \quad\Longleftrightarrow\quad
 a\leq a',\quad U'\subseteq U,\quad c\leq c'.
\]
Given two triples, first dominate their \(a\)- and \(c\)-indices and then use
clusterhood to choose a still later term in the intersection of their
neighborhoods.
Thus \(B\) is directed.
The map \(\phi(a,U,c)=c\) is order preserving, satisfies
\eqref{eq:appE-willard}, and its subnet is eventually in every neighborhood
of \(x\).
\end{proof}

\begin{proposition}\label{prop:appE-divergence-subnet}
If \(x_a\) does not converge to \(x\), there is a Willard subnet no subnet of
which converges to \(x\).
\end{proposition}

\begin{proof}
Some neighborhood \(U\) of \(x\) fails to contain a tail.
Thus the original net is frequently in \(X\setminus U\).
Let \(B\) consist of pairs \((a,c)\) with
\(c\geq a\) and \(x_c\in X\setminus U\), ordered coordinatewise.
Frequency makes \(B\) directed: after dominating two pairs, choose a still
later index whose term lies outside \(U\).  The map
\(\phi(a,c)=c\) is order preserving, satisfies
\eqref{eq:appE-willard}, and produces a subnet entirely in that closed set.
Every further subnet is also entirely there and cannot converge to \(x\).
\end{proof}

\section{Moore--Smith tests for topology}
\label{sec:appE-ms-tests}

\begin{theorem}\label{thm:appE-net-tests}
Let \(X,Y\) be topological spaces.
\begin{enumerate}
\item A point \(x\) lies in \(\overline S\) if and only if a net in \(S\)
      converges to \(x\).
\item A subset \(F\subset X\) is closed if and only if it contains the limit
      of every convergent net in \(F\).
\item A map \(f:X\to Y\) is continuous if and only if it preserves limits of
      all nets.
\item A net in \(\prod_{i\in I}X_i\) converges if and only if every
      coordinate net converges.
\end{enumerate}
\par
\end{theorem}

\begin{proof}
If \(x\in\overline S\), use the neighborhood net and choose
\(x_U\in U\cap S\).
The converse is immediate.
This proves the closed-set criterion.

Continuity preserves eventual membership in inverse images of
neighborhoods.
Conversely, if \(f\) is not continuous at \(x\), some neighborhood \(V\) of
\(f(x)\) has \(x\) in the closure of
\(X\setminus f^{-1}(V)\).
The first assertion gives a net converging to \(x\) whose image never enters
\(V\).

For products, necessity follows from continuity of projections.
A basic product neighborhood restricts only finitely many coordinates.
Directedness supplies one index beyond the finitely many coordinate
thresholds, proving sufficiency.
\end{proof}

\begin{warning}{Sequential shortcuts}
The closure criterion with sequences characterizes sequential spaces, not
all spaces.
Likewise, sequential compactness does not imply compactness in arbitrary
topological spaces.
\Cref{thm:appE-net-tests} has no countability hypothesis because it
uses nets.
\end{warning}

\section{Compactness}
\label{sec:appE-compactness}

\begin{theorem}\label{thm:appE-net-compactness}
A topological space \(K\) is compact if and only if every net in \(K\) has a
convergent Willard subnet.
Equivalently, every net has a convergent order-preserving Moore--Smith
subnet.
If \(K\) is Hausdorff, the subnet limit is unique once the subnet converges.
\end{theorem}

\begin{proof}
For a net \(x_a\), set
\[
 F_a=\overline{\{x_b:b\geq a\}}.
\]
The family \((F_a)\) has the finite-intersection property.
If \(K\) is compact, choose \(x\in\bigcap_aF_a\).
Then \(x_a\) is frequently in every neighborhood of \(x\), so
\Cref{thm:appE-cluster-subnet} gives a convergent subnet.

Conversely, suppose an open cover \(\mathcal U\) has no finite subcover.
Direct its finite subfamilies by inclusion and choose
\[
 x_{\mathcal V}\notin\bigcup_{U\in\mathcal V}U.
\]
If a subnet converged to \(x\), choose \(U\in\mathcal U\) containing \(x\).
It would eventually lie in \(U\).
But \eqref{eq:appE-willard} also forces its indices eventually to contain
\(U\), hence its terms to lie outside \(U\), a contradiction.

If a net converges to two distinct points in a Hausdorff space, it is
eventually in two disjoint neighborhoods, which is impossible.
\end{proof}

\begin{corollary}\label{cor:appE-compact-Hausdorff}
A continuous bijection from a compact space to a Hausdorff space is a
homeomorphism.
\end{corollary}

\begin{proof}
A closed subset of the source is compact, its image is compact, and compact
subsets of a Hausdorff space are closed.
Thus the inverse map is continuous.
\end{proof}

\begin{theorem}\label{thm:appE-dini}
Let \(K\) be compact and let \(f_a\) be a decreasing net of continuous
real-valued functions.
If \(f_a\searrow f\) pointwise and \(f\) is continuous, then
\(f_a\to f\) uniformly.
\end{theorem}

\begin{proof}
For \(\varepsilon>0\), the open sets
\[
                  U_a=\{x:f_a(x)<f(x)+\varepsilon\}
\]
increase with \(a\) and cover \(K\).
A finite subcover is dominated by one index \(a_0\).
For \(a\geq a_0\),
\[
                         0\leq f_a-f<\varepsilon
\]
on \(K\).
\end{proof}

\section{Upper-semicontinuous regularization}
\label{sec:appE-usc}

For \(f:X\to\overline{\mathbf R}\), define
\[
 f^*(x)=\inf_{U\ni x}\sup_{y\in U}f(y).
\tag{E.4}\label{eq:appE-usc-regularization}
\]

\begin{proposition}\label{prop:appE-usc-majorant}
The function \(f^*\) is upper semicontinuous, satisfies \(f\leq f^*\), and
is the smallest upper-semicontinuous majorant of \(f\).
\end{proposition}

\begin{proof}
The condition \(f^*(x)<c\) means that
\(\sup_Uf<c\) for some neighborhood \(U\) of \(x\); it persists throughout
\(U\).
Thus \(\{f^*<c\}\) is open.
Every neighborhood of \(x\) contains \(x\), so \(f\leq f^*\).
If \(g\) is usc and \(f\leq g\), then for \(c>g(x)\) there is a neighborhood
on which \(g<c\), giving \(f^*(x)\leq c\).
Let \(c\downarrow g(x)\).
\end{proof}

\begin{lemma}\label{lem:appE-usc-calculus}
\begin{enumerate}
\item The infimum of any nonempty family of usc functions is usc.
\item A finite maximum of usc functions is usc.
\item For a family \(f_i\), the usc upper envelope is
      \((\sup_i f_i)^*\); regularization is generally necessary.
\item For \(c\in\mathbf R\) and \(\lambda>0\),
\[
 (\sup_i(f_i+c))^*=(\sup_i f_i)^*+c,\qquad
 (\sup_i\lambda f_i)^*=\lambda(\sup_i f_i)^*.
\tag{E.5}\label{eq:appE-usc-affine}
\]
\end{enumerate}
\par
\end{lemma}

\begin{proof}
For an infimum \(h=\inf_i h_i\),
\(\{h<c\}=\bigcup_i\{h_i<c\}\).
For a finite maximum, the same sublevel set is a finite intersection.
The third assertion is
\Cref{prop:appE-usc-majorant}.
Finite translation and positive scaling commute with neighborhood suprema
and infima in \eqref{eq:appE-usc-regularization}.
\end{proof}

\begin{warning}{Regularization is not continuity}
On \([-1,1]\), the function equal to one at \(0\) and zero elsewhere is
upper semicontinuous but discontinuous.
Upper-semicontinuous regularization preserves that spike.
It is an order-theoretic closure, not smoothing.
\end{warning}

\section{Limsup and liminf of nets}
\label{sec:appE-limsup}

For an extended-real net \(r_a\), set
\[
\begin{aligned}
 \limsup_a r_a&=\inf_{a_0}\sup_{a\geq a_0}r_a,\\
 \liminf_a r_a&=\sup_{a_0}\inf_{a\geq a_0}r_a.
\end{aligned}
\tag{E.6}\label{eq:appE-limsup}
\]
These are extended-real quantities and always satisfy
\(\liminf r_a\leq\limsup r_a\).
For a function net \(f_a\), the pointwise upper limit is given by the same
formula at each point, while the usc upper limit is
\[
 \left(\inf_{a_0}\sup_{a\geq a_0}f_a\right)^*.
\tag{E.7}\label{eq:appE-usc-limsup}
\]

\begin{proposition}\label{prop:appE-real-cluster}
If \(r_a\) takes values in a compact interval of \(\mathbf R\), then
\(\limsup r_a\) and \(\liminf r_a\) are respectively the largest and
smallest limits of convergent subnets.
\end{proposition}

\begin{proof}
Compactness gives subnet cluster values.
Every subnet limit lies between the two tail bounds in
\eqref{eq:appE-limsup}.
For the upper endpoint, the original net is frequently in
\((\limsup r_a-\varepsilon,\,+\infty)\) for every \(\varepsilon>0\).
The cluster-subnet construction, directed simultaneously by tails and
\(\varepsilon\downarrow0\), gives a subnet converging to the limsup.
The liminf case is analogous.
\end{proof}

\section{Choquet selection}
\label{sec:appE-choquet}

\begin{theorem}\label{thm:appE-choquet}
Let \(X\) be second countable and let \(\mathcal F\) be a nonempty family of
functions \(X\to\overline{\mathbf R}\).
There is a countable subfamily \((f_j)\) such that
\[
 \left(\sup_{f\in\mathcal F}f\right)^*
 =
 \left(\sup_jf_j\right)^*.
\tag{E.8}\label{eq:appE-choquet}
\]
If \(\mathcal F\) is stable under finite maxima, \(f_j\) may be chosen
increasing.
\end{theorem}

\begin{proof}
Choose a countable base \(\mathcal B\).
For each \(B\in\mathcal B\), let
\[
 M_B=\sup\{f(x):f\in\mathcal F,\ x\in B\}.
\]
When \(M_B\) is finite, choose for every \(q\geq1\) some \(f_{B,q}\) with
\(\sup_Bf_{B,q}>M_B-q^{-1}\).
When \(M_B=+\infty\), choose \(\sup_Bf_{B,q}>q\).
No choice is needed for \(M_B=-\infty\).
Adjoin one fixed member of \(\mathcal F\).
The selected countable family has the same supremum on every basis element,
so formula \eqref{eq:appE-usc-regularization} gives
\eqref{eq:appE-choquet}.
If finite maxima stay in \(\mathcal F\), replace an enumeration \(g_j\) by
\(f_j=\max_{k\leq j}g_k\).
\end{proof}

\begin{warning}{What was selected}
The pointwise suprema in \eqref{eq:appE-choquet} may differ.
Only their usc regularizations agree.
Second countability cannot be omitted unless a different countable
determining family is available.
\end{warning}

\section{Monotone integration firewalls}
\label{sec:appE-monotone-integration}

\begin{theorem}\label{thm:appE-lsc-monotone}
Let \(K\) be compact Hausdorff, let \(\mu\) be a finite Radon measure, and
let \(f_a:K\to[0,+\infty]\) be an increasing net of lower-semicontinuous
functions.
Then
\[
 \int_K\sup_af_a\,\rd \mu
 =
 \sup_a\int_Kf_a\,\rd \mu.
\tag{E.9}\label{eq:appE-lsc-monotone}
\]
\end{theorem}

\begin{proof}
Let \(f=\sup_af_a\), which is lsc.
Radon regularity on compact Hausdorff spaces gives
\[
 \int_Kf\,\rd \mu
 =
 \sup\left\{\int_Kh\,\rd \mu:
 h\in C(K),\ 0\leq h\leq f\right\}.
\]
Fix such an \(h\) and \(\varepsilon>0\).
The open sets
\[
                         \{f_a>h-\varepsilon\}
\]
cover \(K\).
Compactness and directedness give one \(a\) for which the inequality holds
everywhere.
Thus
\[
 \int_Kf_a\,\rd \mu
 \geq\int_Kh\,\rd \mu-\varepsilon\mu(K).
\]
Take the supremum over \(h\) and let \(\varepsilon\downarrow0\).
\end{proof}

\begin{corollary}\label{cor:appE-usc-monotone}
Let \(f_a:K\to[-\infty,C]\) be a decreasing net of usc functions, with
\(C\in\mathbf R\).
Then
\[
 \int_K\inf_af_a\,\rd \mu
 =
 \inf_a\int_Kf_a\,\rd \mu,
\tag{E.10}\label{eq:appE-usc-monotone}
\]
where both sides may equal \(-\infty\).
\end{corollary}

\begin{proof}
Apply \Cref{thm:appE-lsc-monotone} to
\(C-f_a\), which is an increasing nonnegative lsc net.
\end{proof}

\begin{example}[Failure for measurable nets]
\label{ex:appE-measurable-net-failure}
On \([0,1]\) with Lebesgue measure, direct the finite subsets \(F\) by
inclusion and set \(f_F=\mathbf 1_F\).
Then
\[
                         f_F\uparrow1
\]
pointwise, while \(\int f_F=0\) for every \(F\).
Thus the sequential monotone convergence theorem does not extend to
arbitrary measurable nets.
The functions in this example are not lower semicontinuous, so it does not
contradict \Cref{thm:appE-lsc-monotone}.
\end{example}

\begin{warning}{Fatou is also sequential}
The usual proof of Fatou's lemma invokes countable monotone convergence.
There is no unrestricted Fatou lemma for arbitrary nets of measurable
functions.
In this book a net is integrated only through
\Cref{thm:appE-lsc-monotone}, its usc dual, uniform convergence,
or an explicitly selected sequence.
\end{warning}

\section{Weak convergence of measures}
\label{sec:appE-weak-measures}

For compact Hausdorff \(K\), write \(C(K)\) for the real-valued continuous
functions.  A net of finite Radon measures converges weakly if its integrals
against every \(h\in C(K)\) converge.

\begin{proposition}\label{prop:appE-measure-compactness}
The probability measures \(\mathcal P(K)\) form a compact Hausdorff space in
the weak topology.
If \(\mu_a\to\mu\), then
\[
 \liminf_a\int f\,\rd \mu_a\geq\int f\,\rd \mu
\tag{E.11}\label{eq:appE-portmanteau-lsc}
\]
for every bounded-below lsc \(f\), and the reverse limsup inequality holds
for bounded-above usc functions.
\end{proposition}

\begin{proof}
Embed \(\mathcal P(K)\) in
\[
 \prod_{h\in C(K)}[-\lVert h\rVert_\infty,\lVert h\rVert_\infty]
 \quad\text{by}\quad
 \mu\longmapsto\left(\int_Kh\,\rd\mu\right)_h.
\]
The product is compact by the Tychonoff theorem
\cite[Chapter~5]{Kel75}.  A pointwise limit of these coordinate vectors
defines a positive linear functional \(F\) on \(C(K)\) with \(F(1)=1\).
The Riesz representation theorem
\cite[Theorem~2.14]{Rud87} gives a unique probability measure representing
\(F\).  Hence the displayed image is closed and is compact.  Continuous
functions separate Radon measures, so the induced topology is Hausdorff.

For the inequality, first note that a bounded-below lsc function on a
compact Hausdorff space is the pointwise supremum of its continuous
minorants.  Indeed, if \(a<f(x)\), choose \(b\) with \(a<b<f(x)\).
The open set \(\{f>b\}\) contains \(x\); the Urysohn lemma supplies a
continuous function which equals \(b\) at \(x\), equals a fixed lower
bound off that open set, and is everywhere at most \(f\).

We shall also use the integral form of this approximation:
\[
 \int_K f\,\rd\mu
 =\sup\left\{\int_Kh\,\rd\mu:\ h\in C(K),\ h\leq f\right\}.
\tag{E.12}\label{eq:appE-lsc-integral-minorants}
\]
Here are the details.  Subtract a lower bound \(m\), truncate
\(g=f-m\) at height \(N\), and choose a mesh
\(\delta=N/q\).  For \(1\leq j\leq q\), put
\(U_j=\{\min(g,N)>j\delta\}\).  Inner regularity gives compact
\(K_j\subset U_j\), with
\(\delta\sum_j\mu(U_j\setminus K_j)\) arbitrarily small, and the Urysohn
lemma gives \(0\leq\chi_j\leq1\), equal to one on \(K_j\) and supported
in \(U_j\).  Then
\[
 h=m+\delta\sum_{j=1}^{q}\chi_j\in C(K),
 \qquad h\leq f.
\]
Indeed, if \(r\) of the summands are nonzero and \(j_0\) is their largest
index, then \(r\leq j_0\) and \(g>j_0\delta\).  Moreover,
\[
 \int_K(h-m)\,\rd\mu
 \geq\delta\sum_{j=1}^{q}\mu(K_j),
\]
whereas \(\delta\sum_j\mu(U_j)\) is a lower layer-cake sum for
\(\int_K\min(g,N)\,\rd\mu\), with error at most \(\delta\mu(K)\).
First make the regularity error small, then let \(q\to\infty\) and
\(N\to\infty\).  This proves
\eqref{eq:appE-lsc-integral-minorants}.

If \(h\in C(K)\) and \(h\leq f\), then
\[
 \liminf_a\int_K f\,\rd\mu_a
 \geq \lim_a\int_K h\,\rd\mu_a
 =\int_Kh\,\rd\mu.
\]
Taking the supremum over \(h\) and using
\eqref{eq:appE-lsc-integral-minorants} proves
\eqref{eq:appE-portmanteau-lsc}.  Changing signs gives the usc statement.
\end{proof}

If \(K\) is compact metrizable, \(C(K)\) is separable and the weak topology
on \(\mathcal P(K)\) is metrizable.
Only in this case may compactness arguments be rewritten automatically with
subsequences.

\section{Extended-real ledger}
\label{sec:appE-extended-ledger}

\begin{proposition}\label{prop:appE-safe-identities}
For a finite constant \(c\), \(\lambda>0\), and a nonempty family
\((r_i)\subset\overline{\mathbf R}\),
\[
\begin{aligned}
 \sup_i(r_i+c)&=\sup_ir_i+c,&
 \inf_i(r_i+c)&=\inf_ir_i+c,\\
 \sup_i(\lambda r_i)&=\lambda\sup_ir_i,&
 \inf_i(\lambda r_i)&=\lambda\inf_ir_i,\\
 \sup_i\inf_a r_{i,a}&\leq\inf_a\sup_i r_{i,a}.
\end{aligned}
\tag{E.13}\label{eq:appE-safe-identities}
\]
The first four formulas use the standard conventions
\(c+(\pm\infty)=\pm\infty\).
\end{proposition}

\begin{proof}
Translation and positive scaling are order automorphisms of
\(\overline{\mathbf R}\).
For the last inequality, every \(\inf_a r_{i,a}\) is at most every
\(\sup_i r_{i,a}\); take the indicated supremum and infimum.
\end{proof}

\begin{warning}{Forbidden rewrites}
None of the following is valid without finiteness or integrability:
\[
 \sup_i(f_i-g_i)=\sup_if_i-\inf_ig_i,\qquad
 \int(f-g)=\int f-\int g,
\]
when both terms on the right can be infinite.
Likewise, from a carrier inequality \(f\leq p+C\) one may not form
\(f-p\) at a point where \(f=p=-\infty\).
\end{warning}

\section{Unit tests}
\label{sec:appE-unit-tests}

\begin{enumerate}
\item A constant net converges to its constant value.
      Every Willard subnet continues to converge to that value; uniqueness
      of the limit additionally requires a Hausdorff space.
\item The neighborhood net of
      \Cref{ex:appE-neighborhood-net} detects every closure point,
      even when no sequence does.
\item In a product with an uncountable index set, convergence is still
      coordinatewise because a basic neighborhood restricts finitely many
      coordinates.
\item The finite-subset net
      \Cref{ex:appE-measurable-net-failure} is the standard test
      against an illicit net version of monotone convergence.
\item A decreasing continuous net on a compact space with continuous limit
      must pass the uniform estimate of \Cref{thm:appE-dini}.
\item If \(f_a\searrow f\) through usc functions, \(f\) is usc because an
      arbitrary infimum of usc functions is usc.
      The analogous assertion for increasing nets is false without
      regularization.
\end{enumerate}

\section*{Source note}

The subnet, compactness, and Moore--Smith convergence facts follow
\cite[Chapters~4--5]{Wil70}; see also \cite{Kel75}.
The proofs are included because changing the subnet definition can invalidate
a compactness argument.
The envelope and integration statements are elementary consequences of
compactness, Radon regularity, and the explicitly stated semicontinuity
hypotheses.
